\documentclass[12pt,a4paper]{article}
\usepackage[utf8]{inputenc}
\usepackage[T1]{fontenc}
\usepackage[french]{babel}
\usepackage{amsmath, amssymb}
\usepackage{graphicx}
\usepackage{geometry}
\geometry{margin=2cm}
\usepackage{enumitem}
\usepackage{float}
\usepackage{fancyhdr}
\usepackage{titlesec}
\pagestyle{fancy}

% Nettoyer les styles par défaut
\fancyhf{}
 
% En-têtes
\fancyhead[L]{Fiche sur les démonstrations par récurrence} % à gauche
\fancyhead[R]{Matteo Dubresson} % à droite

% Pieds de page
\fancyfoot[C]{\thepage} % numéro de page centré

%Titre et généralités
\title{Fiche sur les démonstrations\\par récurrence}
\date{17/09/2025}
\author{}

%Espace entre les subsections
\titlespacing*{\subsection}{0pt}{5ex plus 0.5ex minus 0.2ex}{1.5ex plus 0.2ex}



\begin{document}

\maketitle

\subsection*{Reconnaitre une récurrence}
Lorsque c'est mentionné dans l'ennoncé \textit{"Montrer que $\forall n \in \mathbb{N^*}$"}, \textit{"Montrer que la suite vérifie ... pour tout $n$"}, ou bien si jamais c'est un exercice de suite, ou une somme quelconque de nombres, c'est très probable qu'une récurrence soit demandée. Mais ce ne sera pas forcément dit explicitement (\textit{"Montrer par récurrence que"}). Il faut donc garder en tête que L'OBJECTIF de la récurrence est de montrer un résultat à n'importe quel ordre, facilement et rapidement. 


\subsection*{Récurrence simple}
Si une propriété est vraie à l'ordre $0$, et à l'ordre $n+1$, alors elle est vrai $\forall n$. Les étapes sont les suivantes :
\begin{itemize}
    \item Initialisation : on montre la propriété à l'odre $0$ ou $1$
    \item Hérédité : on montre la propriété à l'ordre $n+1$ EN AYANT SUPPOSÉ l'ordre $n$ $\to$ il faut que l'hypothèse de récurrence SOIT UTILISÉE (la plupart du temps), et il faut bien le préciser
    \item Conclusion : phrase bateau "On a montré la propriété à l'ordre $0$ et $n+1$, d'où le résultat suivant : (écrire la propriété) $\forall n \in \mathbb{N}$"

\end{itemize}

\subsection*{Récurrence forte}
Même si une récurrence simple pourrait démontrer un certain résultat, dans certains cas une récurrence forte est plus adaptée. Plutot que de supposer un ordre $n$ quelconque, on suppose ici tout les ordres $k\le n$, pour démontrer $n+1$.\\
Les étapes restent les mêmes que pour une récurrence simple, mais l'hérédité peut nécessiter quelques manipulations. 

\subsection*{Applications}
\begin{enumerate}
    \item Soit $(u_n)_{n \in \mathbb{N^*}}$ la suite définie par $u_0 \in ]0;1]$ et par : $ u_{n+1} = \displaystyle \frac{u_n}{2} + \frac{u_n^2}{4}$. \\
    Montrer que pour tout $n \in \mathbb{N}$, on a $0<u_n\le 1$ 
    \vspace{4mm}
    \item Montrer que 
    \[ \forall n \in \mathbb{N^*}, 1+2+3 + ...+ n= \frac{n(n+1)}{2}\]
    \vspace{4mm}
    \item Montrer que 
    \[ \forall n \in \mathbb{N^*}, 1^2+2^2+3^2 + ...+ n^2= \frac{n(n+1)(2n+1)}{6}\]
    \vspace{4mm}
    \item Soit $(u_n)_{n \in \mathbb{N^*}}$ la suite définie par $u_1 =3$ et pour tout $n\ge 1$ par : $ u_{n+1} = \displaystyle \frac{2}{n} \sum_{k=1}^n u_k $. \\
    Montrer que $\forall n \in \mathbb{N^*}$, $u_n =3n$

\end{enumerate}




\end{document}

